\section{Randomized --- Space Amplify}
%%% generated by the blueprint pipeline from TCSlib.Complexity.Randomized.SpaceAmplify
\subsection*{Randomized --- Space Amplify: Introduction}

Declarations of the Lean module \texttt{TCSlib.Complexity.Randomized.SpaceAmplify}.

\subsection*{Randomized --- Space Amplify: Results}

\begin{theorem}[Input-head movement bound]
\lean{Randomized.Amplify.inputPos_stepWith_le}
\label{Randomized.Amplify.inputPos_stepWith_le}
\leanok
\uses{Complexity.LogProg.moveInputPos_val, Turing.Cfg, Turing.NDTM, Turing.NDTM.stepWith, Turing.moveInputPos}
\difficulty{4}
Let $M$ be a nondeterministic Turing machine with Boolean symbols and $K$ work tapes,
let $b$ be a choice bit, and let $c$ be a configuration of $M$ on a Boolean input. After
one step under $b$, the input-head position is at most one greater than its previous
position.
\end{theorem}

\begin{theorem}[Halting composition with a position bound]
\lean{Randomized.Amplify.halts_comp_pos}
\label{Randomized.Amplify.halts_comp_pos}
\leanok
\uses{Randomized.Amplify.inputPos_stepWith_le, Turing.Cfg, Turing.NDTM.runWith, Turing.NDTM.runWith_append, Turing.NDTM.runWith_cons, Turing.NDTM.runWith_of_halt, Turing.NDTM.stepWith}
\difficulty{5}
Let $N$ and $M$ be nondeterministic Turing machines with Boolean symbols, respectively
with $K$ and $K'$ work tapes and state sets $S$ and $S'$, on a fixed Boolean input. Let
$\varphi$ map configurations of $N$ to configurations of $M$. Suppose that whenever a
configuration $c$ of $N$ is nonhalting, one step of $M$ from $\varphi(c)$ under any
choice bit equals the image under $\varphi$ of the corresponding step of $N$. Suppose
there are natural numbers $m,B$ such that, from the image of any halting configuration
$c$ whose input-head position is at most $B$, every choice word of length $m$ leads $M$
to a halting configuration. Then, for every $n\in\mathbb{N}$ and configuration $c$ of
$N$ whose input-head position plus $n$ is at most $B$, if every choice word of length
$n$ leads $N$ from $c$ to a halting configuration, every choice word of length $n+m$
leads $M$ from $\varphi(c)$ to a halting configuration.
\end{theorem}

\begin{theorem}[Halting under composition of simulations]
\lean{Randomized.Amplify.halts_comp}
\label{Randomized.Amplify.halts_comp}
\leanok
\uses{Randomized.Amplify.halts_comp_pos, Turing.Cfg, Turing.NDTM.runWith, Turing.NDTM.stepWith}
\difficulty{1}
Let $N$ and $M$ be nondeterministic Turing machines with Boolean symbols, respectively
with $K$ and $K'$ work tapes and state sets $S$ and $S'$, on a fixed Boolean input. Let
$\varphi$ map configurations of $N$ to configurations of $M$. Suppose that whenever a
configuration $c$ of $N$ is nonhalting, one step of $M$ from $\varphi(c)$ under any
choice bit equals the image under $\varphi$ of the corresponding step of $N$. Suppose
also that there is an $m\in\mathbb{N}$ such that, from the image of any halting
configuration of $N$, every choice word of length $m$ leads $M$ to a halting
configuration. Then, for every $n\in\mathbb{N}$ and configuration $c$ of $N$, if every
choice word of length $n$ leads $N$ from $c$ to a halting configuration, every choice
word of length $n+m$ leads $M$ from $\varphi(c)$ to a halting configuration.
\end{theorem}

\begin{theorem}[Composition of randomized runs]
\lean{Randomized.Amplify.randProb_comp}
\label{Randomized.Amplify.randProb_comp}
\leanok
\uses{Complexity.LogProg.randProb_runWith_choose, Complexity.LogProg.randProb_runWith_mono, Randomized.randProb, Randomized.randProb_const, Randomized.randProb_succ, Turing.Cfg, Turing.NDTM.runWith, Turing.NDTM.runWith_cons, Turing.NDTM.runWith_of_halt, Turing.NDTM.stepWith}
\difficulty{5}
Let $N$ and $M$ be nondeterministic Turing machines with Boolean symbols and
respectively $K$ and $K'$ work tapes, on a fixed Boolean input $x$, and let $\phi$ map
configurations of $N$ to configurations of $M$. Suppose that whenever a configuration
$c$ of $N$ is not halted, one step of $M$ from $\phi(c)$ under any choice bit equals
$\phi$ applied to the corresponding step of $N$. Suppose also that for some
$m\in\mathbb{N}$, every length-$m$ choice word leaves $M$ halted when started from
$\phi(c)$ whenever $c$ is halted; and that, for every halted $c$, the probability that
$M$ satisfies a decidable predicate $Q$ after $m$ choice bits from $\phi(c)$ equals
$g(o(c))$, where $o$ maps configurations of $N$ to Boolean values and $g$ maps Boolean
values to rationals. Then, for every $n\in\mathbb{N}$ and configuration $c$ of $N$ such
that every length-$n$ choice word leaves $N$ halted, the probability that $M$ satisfies
$Q$ after $n+m$ choice bits from $\phi(c)$ equals the probability that $o$ is true after
$n$ choice bits of $N$ times $g(\mathrm{true})$, plus the probability that $o$ is false
after $n$ choice bits of $N$ times $g(\mathrm{false})$.
\end{theorem}

\begin{theorem}[Reachability under composed runs]
\lean{Randomized.Amplify.reach_comp}
\label{Randomized.Amplify.reach_comp}
\leanok
\uses{Turing.Cfg, Turing.NDTM.runWith, Turing.NDTM.runWith_cons, Turing.NDTM.stepWith}
\difficulty{4}
Let $N$ and $M$ be nondeterministic Turing machines with Boolean symbols and
respectively $K$ and $K'$ work tapes, on a fixed Boolean input $x$, and let $\phi$ map
configurations of $N$ to configurations of $M$. Suppose that whenever a configuration
$c$ of $N$ is not halted, one step of $M$ from $\phi(c)$ under any choice bit equals
$\phi$ applied to the corresponding step of $N$. Let $P$ be a predicate on
configurations of $M$. If $P$ holds at $\phi(N\text{'s run from }c\text{ under }u)$ for
every choice word $u$, and if after any such run that halts, $P$ continues to hold after
every further choice word of $M$ from its mapped configuration, then $P$ holds at the
configuration reached by running $M$ on any choice word $w$ from $\phi(c)$.
\end{theorem}

\begin{definition}[Boolean-list summation]
\lean{Randomized.Amplify.summ}
\label{Randomized.Amplify.summ}
\leanok
Define a function from Boolean lists to optional Booleans by returning no value on the
empty list, returning its sole element on a singleton list, and returning false on every
list of length at least two.
\end{definition}

\begin{definition}[Output aggregation]
\lean{Randomized.Amplify.push}
\label{Randomized.Amplify.push}
\leanok
For optional Boolean outputs
$o_1,o_2\in\operatorname{Option}(\{\mathrm{false},\mathrm{true}\})$, define their
aggregate to be $o_1$ if $o_2$ is absent, to be $o_2$ if $o_1$ is absent, and to be the
present value $\mathrm{false}$ if both are present.
\end{definition}

\begin{theorem}[Summary after appending a Boolean value]
\lean{Randomized.Amplify.summ_append}
\label{Randomized.Amplify.summ_append}
\leanok
\uses{Randomized.Amplify.push, Randomized.Amplify.summ}
\difficulty{4}
For a Boolean list $u$, define its summary to be absent when $u$ is empty, its sole
entry when it has length one, and false when it has length at least two. Combining two
optional Boolean values returns the one present if the other is absent, and returns
false if both are present. The summary of $u$ after appending the optional Boolean value
$e$ (when present) is the combination of the summary of $u$ with $e$.
\end{theorem}

\begin{theorem}[Characterization of a singleton true list]
\lean{Randomized.Amplify.summ_eq_some_true}
\label{Randomized.Amplify.summ_eq_some_true}
\leanok
\uses{Randomized.Amplify.summ}
\difficulty{4}
For every finite Boolean list $u$, the summation function returns `some true` exactly
when $u$ is the one-element list $[\mathrm{true}]$.
\end{theorem}

\begin{definition}[Simulation phases]
\lean{Randomized.Amplify.Phase}
\label{Randomized.Amplify.Phase}
\leanok
For a state type $S$ and a natural number $r$, a simulation phase is one of three forms:
running simulation $j$ in state $q$, with optional output $o$ and counter $c$; rewinding
after simulation $j$ has halted, with counter $c$; or seeking the left end of the input
before simulation $j$, with counter $c$. In each form, $j$ ranges from $0$ to $r$
inclusive, and $c$ ranges from $0$ to $r+1$ inclusive.
\end{definition}

\begin{definition}[Boolean counter increment]
\lean{Randomized.Amplify.bump}
\label{Randomized.Amplify.bump}
\leanok
For a natural number $r$, a counter $c\in\{0,\ldots,r+1\}$, and a Boolean $b$, define
the updated counter to be $\min(c+1,r+1)$ when $b$ is true, and $c$ when $b$ is false.
\end{definition}

\begin{definition}[Block coordinate]
\lean{Randomized.Amplify.blk}
\label{Randomized.Amplify.blk}
\leanok
For natural numbers $r$ and $k$, each $t\in\{0,\ldots,(r+1)k-1\}$ has a unique pair of
coordinates in $\{0,\ldots,r\}\times\{0,\ldots,k-1\}$ under the standard finite product
indexing; define its block coordinate to be the first coordinate.
\end{definition}

\begin{definition}[Within-block coordinate]
\lean{Randomized.Amplify.idx}
\label{Randomized.Amplify.idx}
\leanok
For natural numbers $r$ and $k$, each $t\in\{0,\ldots,(r+1)k-1\}$ has a unique pair of
coordinates in $\{0,\ldots,r\}\times\{0,\ldots,k-1\}$ under the standard finite product
indexing; define its within-block coordinate to be the second coordinate.
\end{definition}

\begin{theorem}[Block index under product encoding]
\lean{Randomized.Amplify.blk_mk}
\label{Randomized.Amplify.blk_mk}
\leanok
\uses{Randomized.Amplify.blk}
\difficulty{1}
For natural numbers $r,k$, let $j<r+1$ and $i<k$. Under the standard encoding of pairs
as indices below $(r+1)k$, the block coordinate of the index corresponding to $(j,i)$ is
$j$.
\end{theorem}

\begin{theorem}[Second coordinate of a product index]
\lean{Randomized.Amplify.idx_mk}
\label{Randomized.Amplify.idx_mk}
\leanok
\uses{Randomized.Amplify.idx}
\difficulty{1}
For natural numbers $r,k$, an index $j<r+1$, and an index $i<k$, the second coordinate
of the product index corresponding to $(j,i)$ is $i$.
\end{theorem}

\begin{definition}[Idle machine action]
\lean{Randomized.Amplify.idle}
\label{Randomized.Amplify.idle}
\leanok
\uses{Turing.Action}
For a natural number $k$, state type $\mathrm{St}$, input-head direction
$d\in\{-1,0,1\}$, optional Boolean output $o$, and optional successor state $s$, define
an action on $k$ work tapes that moves the input head by $d$, leaves every work tape
unchanged with its head stationary, emits $o$, and has successor state $s$ (halting when
$s$ is absent).
\end{definition}

\begin{definition}[Amplified machine transition]
\lean{Randomized.Amplify.tr}
\label{Randomized.Amplify.tr}
\leanok
\uses{Randomized.Amplify.Phase, Randomized.Amplify.blk, Randomized.Amplify.bump, Randomized.Amplify.idle, Randomized.Amplify.idx, Randomized.Amplify.push, Turing.Action, Turing.FinNDTM}
Let $N$ be a finite nondeterministic Turing machine over the Boolean alphabet, let $r$
be a natural number, let $\mathrm{dec}:\mathbb{N}\to\{\mathrm{false},\mathrm{true}\}$,
and let $b$ be a choice bit. The transition function of the amplified simulation acts on
phases, an optional input symbol, and the work-tape contents as follows. In a run phase,
it applies $N$'s transition for choice $b$ to the selected simulated run's state and
work tapes, performs that action on the corresponding block of work tapes, and emits no
output; if the simulated transition continues, it updates the accumulated output using
the stated aggregation rule and remains in the run phase, while if it halts, it enters
the rewind phase and increments the counter exactly when the accumulated output is true.
In a rewind phase, if another simulation remains, it moves the input head left and
enters the seek phase for that simulation; otherwise it stays still, outputs
$\mathrm{dec}(c)$, and halts. In a seek phase, on a blank input cell it moves right and
starts the indicated simulation in $N$'s initial state with no accumulated output; on a
present input symbol it moves left and remains in the seek phase.
\end{definition}

\begin{definition}[Amplified finite machine]
\lean{Randomized.Amplify.amp}
\label{Randomized.Amplify.amp}
\leanok
\uses{Randomized.Amplify.Phase, Randomized.Amplify.tr, Turing.FinNDTM}
Given a finite nondeterministic Turing machine $N$ over Boolean symbols, a natural
number $r$, and a function $\mathrm{dec} : \mathbb{N} \to
\{\mathrm{false},\mathrm{true}\}$, define a finite nondeterministic Turing machine with
$(r+1)k$ work tapes, where $k$ is the number of work tapes of $N$, and state space
$\mathrm{Phase}(\mathrm{State}(N),r)$. Its initial state is the run phase indexed by
$0$, initialized with the initial state of $N$ and the additional values `none` and $0$;
its transitions are those specified by the amplification construction from $N$, $r$, and
$\mathrm{dec}$.
\end{definition}

\begin{definition}[Amplified-machine configuration]
\lean{Randomized.Amplify.ACfg}
\label{Randomized.Amplify.ACfg}
\leanok
\uses{Randomized.Amplify.Phase, Turing.Cfg, Turing.FinNDTM}
Given a finite nondeterministic machine $N$ with $k$ work tapes, a repetition count $r$,
and a Boolean input string $x$, an amplified configuration is a configuration on the
input $x$ with $(r+1)k$ work tapes and phase states of the form run, rewind, or seek,
indexed by run numbers in $	ext{Fin}(r+1)$ and counters in $	ext{Fin}(r+2)$.
\end{definition}

\begin{definition}[Embedding a machine configuration]
\lean{Randomized.Amplify.emb}
\label{Randomized.Amplify.emb}
\leanok
\uses{Randomized.Amplify.ACfg, Randomized.Amplify.blk, Randomized.Amplify.bump, Randomized.Amplify.idx, Randomized.Amplify.summ, Turing.Cfg}
Given an amplified configuration $D$, a run index $j\in\text{Fin}(r+1)$, a counter
$c\in\text{Fin}(r+2)$, and a configuration $u$ of $N$ on input $x$, the resulting
amplified configuration has the same input-head position as $u$ and empty output. On
work tapes belonging to run $j$, it copies the tape contents and head positions from
$u$; on all other work tapes, it preserves those of $D$. Its phase state records the
state of $u$ and the summary of its output when $u$ has not halted, and records a rewind
phase with the counter adjusted according to whether that summary is true when $u$ has
halted.
\end{definition}

\begin{theorem}[Step commutes with the configuration embedding]
\lean{Randomized.Amplify.step_emb}
\label{Randomized.Amplify.step_emb}
\leanok
\uses{Randomized.Amplify.ACfg, Randomized.Amplify.amp, Randomized.Amplify.blk, Randomized.Amplify.blk_mk, Randomized.Amplify.emb, Randomized.Amplify.idx_mk, Randomized.Amplify.summ, Randomized.Amplify.summ_append, Randomized.Amplify.tr, Turing.Action.apply, Turing.Cfg, Turing.Cfg.inputSymbol, Turing.Cfg.workTapeSymbols, Turing.NDTM.stepWith}
\difficulty{5}
Let $N$ be a finite nondeterministic Turing machine with Boolean symbols, let $r$ be a
natural number, and let $d\colon\mathbb{N}\to\{0,1\}$. Fix an input $x$, an
amplified-machine configuration $D$, an index $j\in\{0,\ldots,r\}$, a counter
$c\in\{0,\ldots,r+1\}$, a configuration $u$ of $N$ on $x$ whose state is not halted, and
a choice bit $b$. Applying one step with choice $b$ to the amplified machine's embedding
of $u$ at $(j,c)$ gives the same configuration as embedding the result of applying one
step with choice $b$ to $u$ at $(j,c)$.
\end{theorem}

\begin{definition}[Fresh work tapes]
\lean{Randomized.Amplify.Fresh}
\label{Randomized.Amplify.Fresh}
\leanok
\uses{Randomized.Amplify.ACfg, Randomized.Amplify.blk}
Given an amplified configuration $D$ and a natural number $j$, every work tape whose
assigned block index is less than $j$ is blank at every cell and has its head at
position zero.
\end{definition}

\begin{definition}[Rewind configuration condition]
\lean{Randomized.Amplify.RewindAt}
\label{Randomized.Amplify.RewindAt}
\leanok
\uses{Randomized.Amplify.ACfg, Randomized.Amplify.Fresh}
An amplified configuration $C$ is at rewind for run index $j\in\text{Fin}(r+1)$ and
counter $c\in\text{Fin}(r+2)$ precisely when its phase state is that rewind state, its
output is empty, and every work tape whose assigned block index is less than $j$ is
blank at every cell with its head at position zero.
\end{definition}

\begin{theorem}[Embedding a halted run at rewind]
\lean{Randomized.Amplify.rewindAt_emb}
\label{Randomized.Amplify.rewindAt_emb}
\leanok
\uses{Randomized.Amplify.ACfg, Randomized.Amplify.Fresh, Randomized.Amplify.RewindAt, Randomized.Amplify.blk, Randomized.Amplify.bump, Randomized.Amplify.emb, Randomized.Amplify.summ, Turing.Cfg}
\difficulty{4}
Let $N$ be a finite nondeterministic Turing machine over Boolean inputs, let $r$ be a
natural number, and let $x$ be a Boolean input. Take an augmented configuration $D$, an
index $j$ with $0\le j\le r$, a counter $c$ with $0\le c\le r+1$, and a configuration
$u$ of $N$ on $x$. If $u$ is halted and every tape block of $D$ with index greater than
$j$ is blank with its head at cell $0$, then embedding $u$ into block $j$ at counter $c$
produces an augmented configuration in rewind phase at index $j$. Its counter is
increased by one, capped at $r+1$, exactly when the summary of $u$'s output is true; its
output is empty, and every block with index greater than $j$ remains blank with its head
at cell $0$.
\end{theorem}

\begin{definition}[Restarting a run]
\lean{Randomized.Amplify.restart}
\label{Randomized.Amplify.restart}
\leanok
\uses{Randomized.Amplify.ACfg}
Given an amplified configuration $C$, a run index $j\in\text{Fin}(r+1)$, and a counter
$c\in\text{Fin}(r+2)$, restarting produces a configuration that preserves all fields of
$C$ except that its phase state becomes run $j$ in the initial state of $N$, with no
output summary recorded and counter $c$, and its input-head position is set to $1$.
\end{definition}

\begin{theorem}[Restart configuration as an embedding]
\lean{Randomized.Amplify.restart_eq_emb}
\label{Randomized.Amplify.restart_eq_emb}
\leanok
\uses{Randomized.Amplify.ACfg, Randomized.Amplify.blk, Randomized.Amplify.emb, Randomized.Amplify.restart, Randomized.Amplify.summ, Turing.NDTM.initCfg}
\difficulty{4}
Let $N$ be a finite nondeterministic Turing machine with Boolean symbols, let
$r\in\mathbb{N}$, let $x$ be a Boolean input word, and let $C$ be an amplified
configuration. Fix a run index $j<r+1$ and a counter $c<r+2$. Suppose that every work
tape in block $j$ is blank at every cell and has its head at cell $0$, and that the
output is empty. Then restarting $C$ at run $j$ with counter $c$ gives the same
configuration as embedding the restarted configuration at run $j$ with counter $c$,
using $N$'s initial configuration on $x$.
\end{theorem}

\begin{theorem}[Initial configuration as an embedding]
\lean{Randomized.Amplify.init_eq_emb}
\label{Randomized.Amplify.init_eq_emb}
\leanok
\uses{Randomized.Amplify.amp, Randomized.Amplify.emb, Randomized.Amplify.summ, Turing.NDTM.initCfg}
\difficulty{4}
Let $N$ be a finite nondeterministic Turing machine with Boolean symbols, let
$r\in\mathbb{N}$, let $\mathrm{dec}:\mathbb{N}\to\{\mathrm{false},\mathrm{true}\}$, and
let $x$ be a Boolean input word. The initial configuration of the amplified machine on
$x$ equals the embedding of that same configuration into run $0$ with counter $0$, using
the initial configuration of $N$ on $x$.
\end{theorem}

\begin{theorem}[Effect of an idle action]
\lean{Randomized.Amplify.idle_apply}
\label{Randomized.Amplify.idle_apply}
\leanok
\uses{Randomized.Amplify.ACfg, Randomized.Amplify.Phase, Randomized.Amplify.idle, Turing.Action, Turing.Action.apply, Turing.moveInputPos}
\difficulty{3}
Let $N$ be a finite nondeterministic Turing machine with Boolean symbols, let
$r\in\mathbb{N}$, let $\mathrm{dec}:\mathbb{N}\to\{\mathrm{false},\mathrm{true}\}$, and
let $x$ be a finite Boolean input. For any amplified configuration $C$, input-head
movement $\mathrm{inp}\in\{-1,0,1\}$, optional Boolean output $\mathrm{out}$, and
optional successor phase $\mathrm{st}$, applying the action that makes no changes to
work tapes or their head positions, emits $\mathrm{out}$, and sets the state to
$\mathrm{st}$ yields the configuration with state $\mathrm{st}$, input-head position
moved according to $\mathrm{inp}$, unchanged work tapes and work-tape head positions,
and output equal to the prior output followed by $\mathrm{out}$ if present.
\end{theorem}

\begin{theorem}[Step from a specified phase]
\lean{Randomized.Amplify.step_of_state}
\label{Randomized.Amplify.step_of_state}
\leanok
\uses{Randomized.Amplify.ACfg, Randomized.Amplify.Phase, Randomized.Amplify.amp, Randomized.Amplify.tr, Turing.Action.apply, Turing.Cfg.inputSymbol, Turing.Cfg.workTapeSymbols, Turing.NDTM.stepWith}
\difficulty{3}
Let $N$ be a finite-state nondeterministic Turing machine on Boolean inputs, let $r$ be
a natural number, and let $\mathrm{dec}:\mathbb{N}\to\{0,1\}$. If an amplified-machine
configuration $C$ has phase $p$, then for either choice bit $b$, its next configuration
is obtained by applying the amplified transition prescribed for phase $p$, the symbol
under the input head, and the symbols under the work-tape heads.
\end{theorem}

\begin{theorem}[Leftward input-head position bound]
\lean{Randomized.Amplify.moveInputPos_neg_one_le}
\label{Randomized.Amplify.moveInputPos_neg_one_le}
\leanok
\uses{Turing.moveInputPos, Turing.moveInputPos_leftBoundary, Turing.moveInputPos_neg_of_ne_left}
\difficulty{4}
For every $n\in\mathbb{N}$ and every input-head position $p$ among the $n+2$ positions,
moving the head one position to the left with the clamped movement rule gives a position
at most $n$.
\end{theorem}

\begin{theorem}[Restart from the left boundary]
\lean{Randomized.Amplify.step_seek_zero}
\label{Randomized.Amplify.step_seek_zero}
\leanok
\uses{Randomized.Amplify.ACfg, Randomized.Amplify.amp, Randomized.Amplify.idle_apply, Randomized.Amplify.restart, Randomized.Amplify.step_of_state, Randomized.Amplify.tr, Turing.Cfg.inputSymbol, Turing.NDTM.stepWith, Turing.moveInputPos}
\difficulty{5}
Let $N$ be a finite-state nondeterministic Turing machine with Boolean tape symbols, let
$r$ be a natural number, let $\mathrm{dec}$ be a function from natural numbers to
Boolean values, and let $x$ be a finite Boolean input. For any amplified-machine
configuration $C$, run index $j$ among $0,\ldots,r$, counter $c$ among $0,\ldots,r+1$,
and choice bit $b$, if $C$ is in the seek phase for run $j$ with counter $c$ and its
input head is at position $0$, then one step under $b$ gives the configuration obtained
from $C$ by setting its state to the start of run $j$ with counter $c$ and its
input-head position to $1$.
\end{theorem}

\begin{theorem}[Input head moves left during seeking]
\lean{Randomized.Amplify.step_seek_pos}
\label{Randomized.Amplify.step_seek_pos}
\leanok
\uses{Randomized.Amplify.ACfg, Randomized.Amplify.amp, Randomized.Amplify.idle_apply, Randomized.Amplify.step_of_state, Randomized.Amplify.tr, Turing.Cfg.inputSymbol, Turing.NDTM.stepWith, Turing.moveInputPos_neg_of_ne_left}
\difficulty{5}
Let $N$ be a finite-state nondeterministic Turing machine on Boolean input $x$, let $r$
be a natural number, and let $\mathrm{dec}:\mathbb{N}\to\{0,1\}$. If an
amplified-machine configuration $C$ is in the seek phase for run index
$j\in\{0,\ldots,r\}$ and counter $c\in\{0,\ldots,r+1\}$, and its input-head position is
nonzero and at most $|x|$, then one step under either choice bit changes only the
input-head position, decreasing it by one.
\end{theorem}

\begin{theorem}[Returning from the seek phase]
\lean{Randomized.Amplify.seek_run}
\label{Randomized.Amplify.seek_run}
\leanok
\uses{Randomized.Amplify.ACfg, Randomized.Amplify.amp, Randomized.Amplify.restart, Randomized.Amplify.step_seek_pos, Randomized.Amplify.step_seek_zero, Turing.NDTM.runWith, Turing.NDTM.runWith_cons}
\difficulty{5}
Let $N$ be a finite-state nondeterministic Turing machine on Boolean input $x$, let $r$
be a natural number, and let $\mathrm{dec}:\mathbb{N}\to\{0,1\}$. Suppose an
amplified-machine configuration $C$ is in the seek phase for run index
$j\in\{0,\ldots,r\}$ and counter $c\in\{0,\ldots,r+1\}$, with input-head position $p\le
|x|$. Every choice word of length at most $p+1$ leaves all work-tape head positions
unchanged; every choice word of length exactly $p+1$ takes $C$ to the configuration
obtained by restarting run $j$ at counter $c$, with the initial state and input-head
position reset to their starting values while the other configuration data are retained.
\end{theorem}

\begin{theorem}[Restart after a position-dependent rewind]
\lean{Randomized.Amplify.rewind_run_pos}
\label{Randomized.Amplify.rewind_run_pos}
\leanok
\uses{Randomized.Amplify.ACfg, Randomized.Amplify.amp, Randomized.Amplify.idle_apply, Randomized.Amplify.moveInputPos_neg_one_le, Randomized.Amplify.restart, Randomized.Amplify.seek_run, Randomized.Amplify.step_of_state, Randomized.Amplify.tr, Turing.NDTM.runWith, Turing.NDTM.runWith_cons, Turing.NDTM.stepWith, Turing.moveInputPos}
\difficulty{4}
Let $N$ be a finite nondeterministic Turing machine over Boolean inputs, let $r$ be a
natural number, and let $C$ be an augmented configuration in rewind phase at $(j,c)$,
with $0\le j\le r$ and $0\le c\le r+1$. Suppose $j<r$. Then, for every Boolean word $u$
of length at most $q+2$, where $q$ is the input-head position after one attempted move
to the left from $C$, running the amplified machine from $C$ leaves all work-tape head
positions unchanged; and for every such word of length exactly $q+2$, the resulting
configuration is the restart of $C$ at index $j+1$ and counter $c$.
\end{theorem}

\begin{theorem}[Restart after rewinding]
\lean{Randomized.Amplify.rewind_run}
\label{Randomized.Amplify.rewind_run}
\leanok
\uses{Randomized.Amplify.ACfg, Randomized.Amplify.amp, Randomized.Amplify.moveInputPos_neg_one_le, Randomized.Amplify.restart, Randomized.Amplify.rewind_run_pos, Turing.NDTM.runWith}
\difficulty{1}
Let $N$ be a finite nondeterministic Turing machine over Boolean inputs, let $r$ be a
natural number, and let $C$ be an augmented configuration in rewind phase at $(j,c)$,
with $0\le j\le r$ and $0\le c\le r+1$. Suppose $j<r$. Then there is an integer $p$ with
$0\le p\le |x|$ such that, for every Boolean word $u$ of length at most $p+2$, running
the amplified machine from $C$ leaves all work-tape head positions unchanged; and for
every such word of length exactly $p+2$, the resulting configuration is the restart of
$C$ at index $j+1$ and counter $c$.
\end{theorem}

\begin{definition}[Finishing a run]
\lean{Randomized.Amplify.finish}
\label{Randomized.Amplify.finish}
\leanok
\uses{Randomized.Amplify.ACfg}
Given a decoder $d:\mathbb{N}\to\{\mathrm{false},\mathrm{true}\}$, an amplified
configuration $C$, and a counter $c\in\text{Fin}(r+2)$, finishing preserves every field
of $C$ except that the state is set to the halting state and the output is extended by
the single Boolean value $d(c)$ at its end.
\end{definition}

\begin{theorem}[Final rewind step]
\lean{Randomized.Amplify.rewind_last}
\label{Randomized.Amplify.rewind_last}
\leanok
\uses{Randomized.Amplify.ACfg, Randomized.Amplify.amp, Randomized.Amplify.finish, Randomized.Amplify.idle_apply, Randomized.Amplify.step_of_state, Randomized.Amplify.tr, Turing.NDTM.runWith, Turing.NDTM.runWith_cons, Turing.NDTM.runWith_of_halt, Turing.NDTM.stepWith, Turing.moveInputPos_zero}
\difficulty{3}
Let $N$ be a finite nondeterministic Turing machine over Boolean inputs, let $r$ be a
natural number, and let $C$ be an augmented configuration whose phase is rewind at index
$j$ and counter $c$, where $0\le j\le r$ and $0\le c\le r+1$. If $j$ is not less than
$r$, then for every Boolean word $w$, running the amplified machine from $C$ on either
the word formed by prefixing $w$ with true or the word formed by prefixing it with false
yields the same configuration: the configuration obtained from $C$ by halting and
appending the Boolean value of the decision function at $c$ to its output.
\end{theorem}
